\DocumentMetadata{
  pdfversion=2.0,
  pdfstandard=ua-2,
  lang=en-US,
  tagging=on,
  tagging-setup={math/setup=mathml-SE,tabsorder=structure}
}
\documentclass[11pt,letterpaper]{article}
\usepackage[margin=0.68in,includefoot]{geometry}
\usepackage{newcomputermodern}
\usepackage{amsmath,amscd,mathtools}
\usepackage{tikz,adjustbox}
\providecommand{\square}{\mdlgwhtsquare}
\usepackage[hidelinks]{hyperref}
\hypersetup{
  pdftitle={Analysis PhD Exam, January 2023},
  pdfauthor={University of Florida Department of Mathematics},
  pdfsubject={Graduate mathematics examination},
  pdfdisplaydoctitle=true,
  bookmarksopen=true
}
\setcounter{secnumdepth}{0}
\setlength{\parindent}{0pt}
\setlength{\parskip}{0.28em}
\setlength{\emergencystretch}{3em}
\setlength{\leftmargini}{2.15em}
\setlength{\leftmarginii}{2.65em}
\setlength{\leftmarginiii}{3em}
\renewcommand{\labelenumi}{\textbf{\theenumi.}}
\pagestyle{empty}
\raggedbottom
\allowdisplaybreaks
\newenvironment{examproblems}[1]{%
  \begin{enumerate}%
  \setcounter{enumi}{#1}%
  \setlength{\topsep}{0.35em}%
  \setlength{\partopsep}{0pt}%
  \setlength{\itemsep}{0.48em}%
  \setlength{\parsep}{0.18em}%
}{\end{enumerate}}
\newenvironment{examsubparts}{%
  \begin{enumerate}%
  \setlength{\topsep}{0.18em}%
  \setlength{\partopsep}{0pt}%
  \setlength{\itemsep}{0.16em}%
  \setlength{\parsep}{0.1em}%
}{\end{enumerate}}
\newenvironment{examinstructions}{%
  \par\begingroup\itshape
}{%
  \par\endgroup\vspace{0.18em}
}
\makeatletter
\renewcommand\section{\@startsection{section}{1}{0pt}%
  {-1.25ex plus -.25ex minus -.1ex}%
  {0.55ex plus .1ex}%
  {\normalfont\large\bfseries}}
\renewcommand{\@maketitle}{%
  \newpage\null\vskip -1em%
  {\footnotesize\bfseries
    \href{https://www.ufl.edu/}{University of Florida}, \href{https://math.ufl.edu/}{Department of Mathematics}\par}%
  \vskip -0.5em%
  \tagstructbegin{tag=Title}%
    {\LARGE\bfseries\href{https://gma.math.ufl.edu/exams/analysis-phd/}{Analysis PhD Exam}, January 2023\par}%
  \tagstructend%
  \vskip 0.25em%
  \tagmcbegin{artifact}%
    \hrule height 1pt%
  \tagmcend%
  \vskip 1em%
}
\makeatother
\title{Analysis PhD Exam, January 2023}
\author{}
\date{}
\begin{document}
\maketitle
\thispagestyle{empty}
\begin{examinstructions}
Write solutions in a neat and logical fashion, giving complete reasons for all steps. Attempt SIX problems.
\end{examinstructions}
\begin{examproblems}{0}
\item
For any two of the following, either give an example or explain why no example exists.
\begin{examsubparts}
\item[\textnormal{a)}]
a sequence of Lebesgue measurable functions on~\(\mathbb{R}\) which converges in measure but not Lebesgue a.e.
\item[\textnormal{b)}]
a closed set \(E \subset [0,1]\) with positive Lebesgue measure, but which contains no open intervals
\item[\textnormal{c)}]
a decreasing sequence of nonnegative measurable functions on~\(\mathbb{R}\) such that \(\lim \int f_n \ne \int \lim f_n\)
\end{examsubparts}
\item
Prove that for any positive measure~\(\mu\), the space \(L^1(\mu)\) is complete in the \(L^1\) norm.
\item
Let \((\mu_n)\) be a sequence of finite signed measures on a measurable space \((X,\mathcal{M})\). Prove that there exists a finite, positive measure~\(\nu\) on \((X,\mathcal{M})\) such that \(\mu_n \ll \nu\) for all~\(n\).
\item
Suppose that \(f \in L^1[0,1]\) and \(\int_0^x f\,dm=0\) for all \(x \in [0,1]\). What can you conclude about~\(f\)? (Prove your claim.)
\item
Let~\(\mathcal{X}\) be a Banach space. Say that a sequence \((x_n) \subset \mathcal{X}\) converges weakly to \(x \in \mathcal{X}\) if \(f(x_n) \to f(x)\) for all \(f \in \mathcal{X}^*\). Prove the following:
\begin{examsubparts}
\item[\textnormal{a)}]
If \((x_n)\) converges weakly, then \(\sup_n \lVert x_n\rVert < \infty\).
\item[\textnormal{b)}]
Weak limits are unique. (That is, if \(x_n \to x\) weakly and \(x_n \to y\) weakly, then \(x=y\).)
\end{examsubparts}
\item
This problem has two parts.
\begin{examsubparts}
\item[\textnormal{a)}]
State the closed graph theorem.
\item[\textnormal{b)}]
Let~\(\mathcal{X}\) be a Banach space with closed subspaces \(\mathcal{Y},\mathcal{Z} \subset \mathcal{X}\). Suppose that every \(x \in \mathcal{X}\) admits a unique decomposition \(x=y+z\) with \(y \in \mathcal{Y}\), \(z \in \mathcal{Z}\). Prove that there is a constant~\(C\) so that \(\lVert y\rVert \le C\lVert x\rVert\) and \(\lVert z\rVert \le C\lVert x\rVert\).
\end{examsubparts}
\item
Let~\(H\) be a Hilbert space and \(T:H \to H\) a bounded linear operator. Define the adjoint operator \(T^*:H \to H\) and prove that it is bounded, with \(\lVert T^*\rVert=\lVert T\rVert\).
\item
(Short answer, you do not need to give a detailed proof, just a brief explanation.)
\begin{examsubparts}
\item[\textnormal{a)}]
Define the Fourier transform on \(L^1(\mathbb{R})\) and sketch its extension to \(L^2(\mathbb{R})\).
\item[\textnormal{b)}]
Sketch a construction of a bounded linear functional~\(\lambda\) on \(L^\infty(\mathbb{R})\) which is not of the form \(\lambda(f)=\int f(x)g(x)\,dx\) for any \(L^1\) function~\(g\).
\item[\textnormal{c)}]
Sketch a proof of the fact that there is no norm under which \(c_{00}\) is complete.
\end{examsubparts}
\end{examproblems}
\end{document}
